\section{Velocity--score duality and the marginal-preserving SDE family}
\label{appendix:score_velocity}

This appendix establishes the two structural facts on which the model-agnostic method rests: that for an affine Gaussian path the velocity and the score determine each other in closed form (Proposition~\ref{prop:score_velocity}), and that the probability-flow ODE belongs to a one-parameter family of SDEs sharing the same marginals (Proposition~\ref{prop:marginal_preservation}). Both are known \cite{albergo_stochastic_2023, song_score-based_2021, ma_sit_2024}; we state and prove them in the notation of Section~\ref{subsec:flow_models} for completeness.

Throughout, consider the affine Gaussian path of Eq.~\eqref{eq:general_interpolant},
\begin{align}\label{eq:app_path}
    \bx_\tau = \alpha_\tau\bm a_0 + \beta_\tau\bx_1 + \sigpath_\tau\latentnoise,
    \qquad \latentnoise\sim\mathcal{N}(0,\bI)\ \text{independent of}\ \bx_1,
\end{align}
with $\bm a_0$ deterministic, $\beta_\tau>0$ for $\tau\in(0,1]$, and $\sigpath_\tau>0$ for $\tau\in[0,1)$. Let $p_\tau$ denote the marginal density of $\bx_\tau$ and $\genscore_\tau=\nabla_{\bx}\log p_\tau$ its score.

\begin{lemma}[Tweedie identities for the affine path]\label{lem:tweedie}
Writing $\bm m_\tau := \alpha_\tau\bm a_0 + \beta_\tau\bx_1$, so that $\bx_\tau = \bm m_\tau + \sigpath_\tau\latentnoise$ with $\sigpath_\tau\latentnoise$ Gaussian of covariance $\sigpath_\tau^2\bI$ independent of $\bm m_\tau$,
\begin{align}
    \E[\bm m_\tau \mid \bx_\tau] &= \bx_\tau + \sigpath_\tau^2\,\genscore_\tau(\bx_\tau), \label{eq:tweedie_mean}\\
    \E[\latentnoise \mid \bx_\tau] &= -\sigpath_\tau\,\genscore_\tau(\bx_\tau), \label{eq:tweedie_eps}\\
    \E[\bx_1 \mid \bx_\tau] &= \tfrac{1}{\beta_\tau}\!\left(\bx_\tau + \sigpath_\tau^2\genscore_\tau(\bx_\tau) - \alpha_\tau\bm a_0\right). \label{eq:tweedie_x1}
\end{align}
\end{lemma}
\begin{proof}
Equation~\eqref{eq:tweedie_mean} is the Gaussian Tweedie/Miyasawa identity for the additive-Gaussian-noise model $\bx_\tau = \bm m_\tau + \sigpath_\tau\latentnoise$: differentiating $p_\tau(\bx)=\int \mathcal{N}(\bx;\bm m,\sigpath_\tau^2\bI)\,q(\bm m)\,\rd\bm m$ gives $\nabla_{\bx} p_\tau = \int \sigpath_\tau^{-2}(\bm m-\bx)\mathcal{N}(\bx;\bm m,\sigpath_\tau^2\bI)q(\bm m)\rd\bm m$, hence $\sigpath_\tau^2\genscore_\tau(\bx) = \E[\bm m_\tau\mid\bx_\tau]-\bx$, which is \eqref{eq:tweedie_mean}. Since $\latentnoise=(\bx_\tau-\bm m_\tau)/\sigpath_\tau$, \eqref{eq:tweedie_eps} follows. Finally $\bx_1=(\bm m_\tau-\alpha_\tau\bm a_0)/\beta_\tau$ gives \eqref{eq:tweedie_x1}.
\end{proof}

\begin{proposition}[Velocity--score identity]\label{prop:score_velocity}
The marginal velocity \eqref{eq:general_velocity} and the score of the path \eqref{eq:app_path} satisfy the exact linear relation
\begin{align}\label{eq:app_velocity_of_score}
    \fmvel_\tau(\bx)
    = \frac{\dot\beta_\tau}{\beta_\tau}\bx
    + \left(\dot\alpha_\tau - \frac{\dot\beta_\tau\alpha_\tau}{\beta_\tau}\right)\bm a_0
    + \sigpath_\tau\!\left(\frac{\dot\beta_\tau}{\beta_\tau}\sigpath_\tau - \dot\sigpath_\tau\right)\genscore_\tau(\bx),
\end{align}
where the score coefficient is the velocity--score coefficient $\vscoef_\tau=\sigpath_\tau(\tfrac{\dot\beta_\tau}{\beta_\tau}\sigpath_\tau-\dot\sigpath_\tau)$ of Eq.~\eqref{eq:vscoef_def}, equivalently
\begin{align}\label{eq:app_score_of_velocity}
    \genscore_\tau(\bx)
    = \frac{\beta_\tau\,\fmvel_\tau(\bx) - \dot\beta_\tau\,\bx - (\dot\alpha_\tau\beta_\tau - \dot\beta_\tau\alpha_\tau)\,\bm a_0}
           {\sigpath_\tau\,(\dot\beta_\tau\sigpath_\tau - \dot\sigpath_\tau\beta_\tau)} .
\end{align}
In particular, for flow matching ($\bm a_0=\bm 0$, $\sigpath_\tau=\alpha_\tau$) this is Eq.~\eqref{eq:fm_score}, and for the stochastic interpolant the corresponding statement for the SDE drift $\drift_\theta=\fmvel_\tau+\tfrac12\gamma_\tau^2\genscore_\tau$ is Eq.~\eqref{eq:SI_score}.
\end{proposition}
\begin{proof}
By definition $\fmvel_\tau(\bx)=\E[\dot\alpha_\tau\bm a_0 + \dot\beta_\tau\bx_1 + \dot\sigpath_\tau\latentnoise\mid\bx_\tau=\bx]$. Since $\bm a_0$ is deterministic, substitute \eqref{eq:tweedie_x1} and \eqref{eq:tweedie_eps}:
\begin{align}
    \fmvel_\tau(\bx)
    &= \dot\alpha_\tau\bm a_0
    + \frac{\dot\beta_\tau}{\beta_\tau}\!\left(\bx + \sigpath_\tau^2\genscore_\tau - \alpha_\tau\bm a_0\right)
    - \dot\sigpath_\tau\sigpath_\tau\genscore_\tau \nonumber\\
    &= \frac{\dot\beta_\tau}{\beta_\tau}\bx
    + \left(\dot\alpha_\tau - \frac{\dot\beta_\tau\alpha_\tau}{\beta_\tau}\right)\bm a_0
    + \left(\frac{\dot\beta_\tau}{\beta_\tau}\sigpath_\tau^2 - \dot\sigpath_\tau\sigpath_\tau\right)\genscore_\tau,
\end{align}
which is \eqref{eq:app_velocity_of_score}. Multiplying the isolated score term by $\beta_\tau$ and rearranging yields \eqref{eq:app_score_of_velocity}. The flow matching specialization is immediate. For the stochastic interpolant the trained object is the SDE drift rather than the ODE velocity; the analogous score relation \eqref{eq:SI_score} is derived in \cite{chen_probabilistic_2024} and recovered here by adding $\tfrac12\gamma_\tau^2\genscore_\tau$ to $\fmvel_\tau$.
\end{proof}

\begin{proposition}[Marginal-preserving SDE family]\label{prop:marginal_preservation}
Let $p_\tau$ be the marginals of the path \eqref{eq:app_path} with velocity $\fmvel_\tau$ and score $\genscore_\tau$. For any measurable $\gdiff_\tau\ge0$, the SDE
\begin{align}\label{eq:app_sde}
    \rd\bX_\tau = \left[\fmvel_\tau(\bX_\tau) + \tfrac12\gdiff_\tau^2\genscore_\tau(\bX_\tau)\right]\rd\tau + \gdiff_\tau\,\rd\bW_\tau,
    \qquad \bX_0\sim p_0,
\end{align}
has marginal law $p_\tau$ for all $\tau\in[0,1]$. Setting $\gdiff_\tau=0$ gives the probability-flow ODE.
\end{proposition}
\begin{proof}
By construction the marginals satisfy the continuity (transport) equation associated with the velocity,
\begin{align}\label{eq:app_continuity}
    \partial_\tau p_\tau = -\nabla\!\cdot\!\left(\fmvel_\tau\,p_\tau\right),
\end{align}
since $\fmvel_\tau$ is the conditional expectation generating the path \cite{lipman_flow_2023, albergo_stochastic_2023}. The Fokker--Planck equation of \eqref{eq:app_sde} for a density $q_\tau$ reads
\begin{align}
    \partial_\tau q_\tau
    = -\nabla\!\cdot\!\left(\left[\fmvel_\tau + \tfrac12\gdiff_\tau^2\genscore_\tau\right]q_\tau\right)
    + \tfrac12\gdiff_\tau^2\Delta q_\tau .
\end{align}
Evaluating the right-hand side at $q_\tau=p_\tau$ and using $\Delta p_\tau = \nabla\!\cdot\!(\nabla p_\tau) = \nabla\!\cdot\!(p_\tau\nabla\log p_\tau) = \nabla\!\cdot\!(p_\tau\genscore_\tau)$,
\begin{align}
    -\nabla\!\cdot\!\left(\fmvel_\tau p_\tau\right)
    - \tfrac12\gdiff_\tau^2\nabla\!\cdot\!\left(\genscore_\tau p_\tau\right)
    + \tfrac12\gdiff_\tau^2\nabla\!\cdot\!\left(p_\tau\genscore_\tau\right)
    = -\nabla\!\cdot\!\left(\fmvel_\tau p_\tau\right)
    = \partial_\tau p_\tau,
\end{align}
by \eqref{eq:app_continuity}. Hence $p_\tau$ solves the Fokker--Planck equation of \eqref{eq:app_sde}. As the solution with initial datum $p_0$ is unique under standard regularity assumptions, the law of $\bX_\tau$ equals $p_\tau$ for all $\tau$.
\end{proof}

\begin{remark}
Propositions~\ref{prop:score_velocity}--\ref{prop:marginal_preservation} are what make Theorem~\ref{theorem:unified_posterior} (and the path-space Theorem~\ref{theorem:posterior_sde}) applicable to a deterministically trained flow matching model: although the model is sampled as an ODE ($\gdiff_\tau=0$), one may pick any $\gdiff_\tau>0$, recover $\genscore_\tau$ from the velocity by Proposition~\ref{prop:score_velocity}, and run the equal-marginal SDE \eqref{eq:app_sde}. The conditional path of Proposition~\ref{prop:conditional_path} then inherits the same duality, and the uncontrolled SDE has terminal law $\conddist{\bx_1}{\bx_0}$, the only hypothesis the posterior samplers require.
\end{remark}
